\import{algebra/monoid.tex}
\import{set/finite-fold.tex}

\section{Enumeration-independent finite folds}

\begin{proposition}\label{spike_sequence_append_range}
    Suppose $s$ is a finite sequence in $A$ and $x\in A$.
    Then
        $\ran{\sequenceAppend{s}{x}}=\cons{x}{\ran{s}}$.
\end{proposition}
\begin{proof}
    Follows by \auto{sequence_append,ran_cons}.
\end{proof}

\begin{proposition}\label{spike_sequence_append_injective_elim}
    Suppose $s$ is a finite sequence in $A$ and $x\in A$.
    Suppose $\sequenceAppend{s}{x}$ is injective.
    Then $s$ is injective and $x\notin\ran{s}$.
\end{proposition}
\begin{proof}
    Let $t=\sequenceAppend{s}{x}$.
    $t$ is injective by \auto.
    $s$ is a function and $t$ is a function by
        \auto{finite_sequence_domain,funs_is_function,
            sequence_append_finite_sequence}.
    $s\subseteq t$ by \auto{spike_sequence_append_range,sequence_append,
        subseteq,cons_iff}.
    Show that $s$ is injective.
    \begin{subproof}
        Show that for all $i,j,u$ such that $(i,u),(j,u)\in s$ we have
            $i=j$.
        \begin{subproof}
            Fix $i,j,u$.
            Assume $(i,u),(j,u)\in s$.
            Hence $(i,u),(j,u)\in t$ by \auto{subseteq}.
            Follows by \auto{injective}.
        \end{subproof}
        Follows by \auto{injective}.
    \end{subproof}
    Show that $x\notin\ran{s}$.
    \begin{subproof}
        Suppose not.
        Take $i\in\dom{s}$ such that $s(i)=x$ by
            \auto{fun_ran_iff}.
        $\dom{s}\notin\dom{s}$ by \auto{in_irrefl}.
        $i\in\dom{t}$ by
            \auto{sequence_append_domain,subseteq_self_suc_intro,subseteq}.
        $\dom{s}\in\dom{t}$ by
            \auto{sequence_append_domain,suc_intro_self}.
        $t(i)=s(i)$ by \auto{sequence_append_old_value}.
        $t(\dom{s})=x$ by \auto{sequence_append_last_value}.
        Hence $i=\dom{s}$ by \auto{injective_function}.
        Contradiction.
    \end{subproof}
    Follows.
\end{proof}

\begin{proposition}\label{spike_sequence_append_injective_intro}
    Suppose $s$ is a finite sequence in $A$ and $x\in A$.
    Suppose $s$ is injective and $x\notin\ran{s}$.
    Then $\sequenceAppend{s}{x}$ is injective.
\end{proposition}
\begin{proof}
    Let $t=\sequenceAppend{s}{x}$.
    $s$ is a function and $t$ is a function by
        \auto{finite_sequence_domain,funs_is_function,
            sequence_append_finite_sequence}.
    Show that for all $i,j\in\dom{t}$ such that $t(i)=t(j)$ we have
        $i=j$.
    \begin{subproof}
        Fix $i,j\in\dom{t}$.
        Assume $t(i)=t(j)$.
        $i\in\dom{s}$ or $i=\dom{s}$ by
            \auto{sequence_append_domain,suc_iff}.
        $j\in\dom{s}$ or $j=\dom{s}$ by
            \auto{sequence_append_domain,suc_iff}.
        \begin{byCase}
        \caseOf{$i\in\dom{s}$ and $j\in\dom{s}$.}
            $t(i)=s(i)$ and $t(j)=s(j)$ by
                \auto{sequence_append_old_value}.
            Hence $s(i)=s(j)$ by \auto.
            Follows by \auto{injective_function}.
        \caseOf{It is not the case that
            $i\in\dom{s}$ and $j\in\dom{s}$.}
            $i=\dom{s}$ or $j=\dom{s}$ by \auto.
            \begin{byCase}
            \caseOf{$i=\dom{s}$.}
                $t(i)=x$ by \auto{sequence_append_last_value}.
                Show that $j=\dom{s}$.
                \begin{subproof}
                    Suppose not.
                    Hence $j\in\dom{s}$ by \auto.
                    $t(j)=s(j)$ by \auto{sequence_append_old_value}.
                    Hence $s(j)=x$ by \auto.
                    Thus $x\in\ran{s}$ by \auto{fun_ran_iff}.
                    Contradiction.
                \end{subproof}
                Follows by \auto.
            \caseOf{$i\neq\dom{s}$.}
                Hence $j=\dom{s}$ by \auto.
                $t(j)=x$ by \auto{sequence_append_last_value}.
                Show that $i=\dom{s}$.
                \begin{subproof}
                    Suppose not.
                    Hence $i\in\dom{s}$ by \auto.
                    $t(i)=s(i)$ by \auto{sequence_append_old_value}.
                    Hence $s(i)=x$ by \auto.
                    Thus $x\in\ran{s}$ by \auto{fun_ran_iff}.
                    Contradiction.
                \end{subproof}
                Follows.
            \end{byCase}
        \end{byCase}
    \end{subproof}
    Follows by \auto{injective_function}.
\end{proof}

\begin{proposition}\label{spike_map_finite_sequence}
    Suppose $s$ is a finite sequence in $A$ and $f\in\funs{A}{B}$.
    Then $f\circ s$ is a finite sequence in $B$.
\end{proposition}
\begin{proof}
    $s\in\funs{\dom{s}}{A}$ and $\dom{s}$ is a natural number by
        \auto{finite_sequence_domain}.
    Hence $f\circ s\in\funs{\dom{s}}{B}$ by \auto{funs_circ}.
    Follows by \auto{finite_sequence}.
\end{proof}

\begin{proposition}\label{spike_map_sequence_append}
    Suppose $s$ is a finite sequence in $A$.
    Suppose $f\in\funs{A}{B}$ and $x\in A$.
    Then
        $f\circ\sequenceAppend{s}{x}
            =\sequenceAppend{f\circ s}{f(x)}$.
\end{proposition}
\begin{proof}
    Let $t=\sequenceAppend{s}{x}$.
    Let $u=f\circ t$ and $v=\sequenceAppend{f\circ s}{f(x)}$.
    $s\in\funs{\dom{s}}{A}$ by \auto{finite_sequence_domain}.
    $f\circ s$ is a finite sequence in $B$ by
        \auto{spike_map_finite_sequence}.
    $f(x)\in B$ by \auto{funs_type_apply}.
    $t$ is a finite sequence in $A$ by
        \auto{sequence_append_finite_sequence}.
    Hence $t\in\funs{\dom{t}}{A}$ by
        \auto{finite_sequence_domain}.
    $f$ is composable with $t$ and $f$ is composable with $s$ by
        \auto{funs_ran,funs}.
    Hence $u$ is a finite sequence in $B$ by
        \auto{spike_map_finite_sequence}.
    $v$ is a finite sequence in $B$ by
        \auto{sequence_append_finite_sequence}.
    Hence $u$ is a function by
        \auto{finite_sequence_domain,funs_is_function}.
    $v$ is a function by
        \auto{finite_sequence_domain,funs_is_function}.
    $\dom{u}=\dom{t}=\suc{\dom{s}}$ by
        \auto{funs_circ,finite_sequence_domain,funs_elim,
            sequence_append_domain}.
    $\dom{v}=\suc{\dom{f\circ s}}=\suc{\dom{s}}$ by
        \auto{sequence_append_domain,funs_circ,
            finite_sequence_domain,funs_elim}.
    Hence $\dom{u}=\dom{v}$ by \auto.
    Show that for every $i\in\dom{u}$ we have $u(i)=v(i)$.
    \begin{subproof}
        Fix $i\in\dom{u}$.
        Hence $i\in\suc{\dom{s}}$ by \auto.
        $i\in\dom{s}$ or $i=\dom{s}$ by \auto{suc_iff}.
        \begin{byCase}
        \caseOf{$i\in\dom{s}$.}
            Hence $i\in\dom{t}$ by
                \auto{sequence_append_domain,
                    subseteq_self_suc_intro,subseteq}.
            $(f\circ t)(i)=f(t(i))$ by
                \auto{circ_apply,funs_is_function}.
            Hence $u(i)=f(t(i))$ by \auto.
            $t(i)=s(i)$ by \auto{sequence_append_old_value}.
            $i\in\dom{f\circ s}$ by
                \auto{funs_circ,finite_sequence_domain,funs_elim}.
            $v(i)=(f\circ s)(i)$ by
                \auto{sequence_append_old_value}.
            $(f\circ s)(i)=f(s(i))$ by
                \auto{circ_apply,funs_is_function}.
            Follows by \auto.
        \caseOf{$i\notin\dom{s}$.}
            Hence $i=\dom{s}$ by \auto.
            $i\in\dom{t}$ by
                \auto{sequence_append_domain,suc_intro_self}.
            $(f\circ t)(i)=f(t(i))$ by
                \auto{circ_apply,funs_is_function}.
            Hence $u(i)=f(t(i))$ by \auto.
            $t(i)=x$ by \auto{sequence_append_last_value}.
            $\dom{f\circ s}=\dom{s}$ by
                \auto{funs_circ,finite_sequence_domain,funs_elim}.
            Hence $v(i)=f(x)$ by
                \auto{sequence_append_last_value}.
            Follows by \auto.
        \end{byCase}
    \end{subproof}
    Follows by \auto{funext_on_domain}.
\end{proof}

\begin{proposition}\label{spike_mapped_append_fold}
    Let $M$ be a monoid.
    Suppose $s$ is a finite sequence in $A$.
    Suppose $f\in\funs{A}{\carrier[M]}$ and $x\in A$.
    Then
        $\sequenceFold{\mul[M]}{\neutral[M]}{
            f\circ\sequenceAppend{s}{x}}
        =\mul[M]
            (\sequenceFold{\mul[M]}{\neutral[M]}{f\circ s},f(x))$.
\end{proposition}
\begin{proof}
    $f\circ s$ is a finite sequence in $\carrier[M]$ by
        \auto{spike_map_finite_sequence}.
    $f(x)\in\carrier[M]$ by \auto{funs_type_apply}.
    Follows by
        \auto{spike_map_sequence_append,sequence_fold_append}.
\end{proof}

\begin{proposition}\label{spike_mapped_fold_type}
    Let $M$ be a monoid.
    Suppose $s$ is a finite sequence in $A$.
    Suppose $f\in\funs{A}{\carrier[M]}$.
    Then
        $\sequenceFold{\mul[M]}{\neutral[M]}{f\circ s}
            \in\carrier[M]$.
\end{proposition}
\begin{proof}
    $f\circ s$ is a finite sequence in $\carrier[M]$ by
        \auto{spike_map_finite_sequence}.
    $\neutral[M]\in\carrier[M]$ by
        \auto{monoid_is_unital_magma,unitalmagma_type}.
    $\mul[M]\in
        \funs{\carrier[M]\times\carrier[M]}{\carrier[M]}$ by
        \auto{monoid_is_magma,magma_mul_graph}.
    Follows by \auto{sequence_fold_typed}.
\end{proof}

\begin{proposition}\label{spike_commutative_monoid_exchange}
    Let $M$ be a monoid.
    Suppose for all $x,y\in\carrier[M]$ we have
        $\mul[M](x,y)=\mul[M](y,x)$.
    Let $p,a,b\in\carrier[M]$.
    Then
        $\mul[M](\mul[M](p,a),b)
            =\mul[M](\mul[M](p,b),a)$.
\end{proposition}
\begin{proof}
    $\mul[M](a,b)=\mul[M](b,a)$ by \auto.
    $\mul[M](p,\mul[M](a,b))
        =\mul[M](\mul[M](p,a),b)$ by \auto{monoid_assoc}.
    $\mul[M](p,\mul[M](b,a))
        =\mul[M](\mul[M](p,b),a)$ by \auto{monoid_assoc}.
    Follows by \auto.
\end{proof}

\begin{theorem}\label{spike_fold_extract_value}
    Let $M$ be a monoid.
    Suppose for all $u,v\in\carrier[M]$ we have
        $\mul[M](u,v)=\mul[M](v,u)$.
    Suppose $s$ is a finite sequence in $A$ and $s$ is injective.
    Suppose $f\in\funs{A}{\carrier[M]}$ and $x\in\ran{s}$.
    Then there exists $r$ such that
        $r$ is a finite sequence in $A$ and
        $r$ is injective and
        $\ran{r}=\ran{s}\setminus\{x\}$ and
        $\sequenceFold{\mul[M]}{\neutral[M]}{f\circ s}
            =\mul[M]
                (\sequenceFold{\mul[M]}{\neutral[M]}{f\circ r},
                 f(x))$.
\end{theorem}
\begin{proof}
    Let $C=\{t\in\finiteSequences{A}\mid
        \text{if $t$ is injective, then for every $a\in\ran{t}$
        there exists $r$ such that
            $r$ is a finite sequence in $A$ and
            $r$ is injective and
            $\ran{r}=\ran{t}\setminus\{a\}$ and
            $\sequenceFold{\mul[M]}{\neutral[M]}{f\circ t}
                =\mul[M]
                    (\sequenceFold{\mul[M]}{\neutral[M]}{f\circ r},
                     f(a))$}\}$.
    $\emptyset\in C$ by
        \auto{finite_sequences_iff,emptyset_finite_sequence,
            ran_emptyset,emptyset}.
    Show that for every $t,y$ such that
        $t\in C$ and $y\in A$ and $t$ is a finite sequence in $A$
        we have $\sequenceAppend{t}{y}\in C$.
    \begin{subproof}
        Fix $t,y$.
        Assume $t\in C$ and $y\in A$ and
            $t$ is a finite sequence in $A$.
        Let $q=\sequenceAppend{t}{y}$.
        $q$ is a finite sequence in $A$ by
            \auto{sequence_append_finite_sequence}.
        Hence $q\in\finiteSequences{A}$ by
            \auto{finite_sequences_iff}.
        Show that if $q$ is injective, then for every $a\in\ran{q}$
            there exists $r$ such that
                $r$ is a finite sequence in $A$ and
                $r$ is injective and
                $\ran{r}=\ran{q}\setminus\{a\}$ and
                $\sequenceFold{\mul[M]}{\neutral[M]}{f\circ q}
                    =\mul[M]
                        (\sequenceFold{\mul[M]}{\neutral[M]}{f\circ r},
                         f(a))$.
        \begin{subproof}
            Assume $q$ is injective.
            Hence $t$ is injective and $y\notin\ran{t}$ by
                \auto{spike_sequence_append_injective_elim}.
            Fix $a\in\ran{q}$.
            $\ran{q}=\cons{y}{\ran{t}}$ by
                \auto{spike_sequence_append_range}.
            \begin{byCase}
            \caseOf{$a=y$.}
                $\ran{t}=\ran{q}\setminus\{a\}$ by
                    \auto{setext,setminus,cons_iff,singleton_iff}.
                $\sequenceFold{\mul[M]}{\neutral[M]}{f\circ q}
                    =\mul[M]
                        (\sequenceFold{\mul[M]}{\neutral[M]}{f\circ t},
                         f(a))$ by
                    \auto{spike_mapped_append_fold}.
                Follows by \auto.
            \caseOf{$a\neq y$.}
                Hence $a\in\ran{t}$ by \auto{cons_iff}.
                Take $r$ such that
                    $r$ is a finite sequence in $A$ and
                    $r$ is injective and
                    $\ran{r}=\ran{t}\setminus\{a\}$ and
                    $\sequenceFold{\mul[M]}{\neutral[M]}{f\circ t}
                        =\mul[M]
                            (\sequenceFold{\mul[M]}{\neutral[M]}{f\circ r},
                             f(a))$ by \auto.
                $y\notin\ran{r}$ by \auto{setminus}.
                Let $w=\sequenceAppend{r}{y}$.
                $w$ is a finite sequence in $A$ by
                    \auto{sequence_append_finite_sequence}.
                $w$ is injective by
                    \auto{spike_sequence_append_injective_intro}.
                $\ran{w}=\cons{y}{\ran{r}}$ by
                    \auto{spike_sequence_append_range}.
                Hence $\ran{w}=\ran{q}\setminus\{a\}$ by
                    \auto{setext,setminus,cons_iff,singleton_iff}.
                Let $p=\sequenceFold{\mul[M]}{\neutral[M]}{f\circ r}$.
                $p\in\carrier[M]$ by
                    \auto{spike_mapped_fold_type}.
                $a\in A$ by
                    \auto{finite_sequence_domain,funs_ran,subseteq}.
                Hence $f(a)\in\carrier[M]$ by
                    \auto{funs_type_apply}.
                $f(y)\in\carrier[M]$ by \auto{funs_type_apply}.
                $\sequenceFold{\mul[M]}{\neutral[M]}{f\circ q}
                    =\mul[M]
                        (\sequenceFold{\mul[M]}{\neutral[M]}{f\circ t},
                         f(y))$ by
                    \auto{spike_mapped_append_fold}.
                Hence
                    $\sequenceFold{\mul[M]}{\neutral[M]}{f\circ q}
                        =\mul[M](\mul[M](p,f(a)),f(y))$ by
                    \auto.
                $\sequenceFold{\mul[M]}{\neutral[M]}{f\circ w}
                    =\mul[M](p,f(y))$ by
                    \auto{spike_mapped_append_fold}.
                $\mul[M](\mul[M](p,f(a)),f(y))
                    =\mul[M](\mul[M](p,f(y)),f(a))$ by
                    \auto{spike_commutative_monoid_exchange}.
                Hence
                    $\sequenceFold{\mul[M]}{\neutral[M]}{f\circ q}
                        =\mul[M]
                            (\sequenceFold{\mul[M]}{\neutral[M]}{f\circ w},
                             f(a))$ by \auto.
                Follows by \auto.
            \end{byCase}
        \end{subproof}
        Follows by \auto.
    \end{subproof}
    $s\in C$ by \auto{finite_sequence_induction}.
    Follows by \auto.
\end{proof}

\begin{theorem}\label{spike_injective_sequence_fold_invariant}
    Let $M$ be a monoid.
    Suppose for all $u,v\in\carrier[M]$ we have
        $\mul[M](u,v)=\mul[M](v,u)$.
    Suppose $s$ is a finite sequence in $A$ and
        $t$ is a finite sequence in $A$.
    Suppose $s$ is injective and $t$ is injective.
    Suppose $\ran{s}=\ran{t}$.
    Suppose $f\in\funs{A}{\carrier[M]}$.
    Then
        $\sequenceFold{\mul[M]}{\neutral[M]}{f\circ s}
            =\sequenceFold{\mul[M]}{\neutral[M]}{f\circ t}$.
\end{theorem}
\begin{proof}
    Let $C=\{q\in\finiteSequences{A}\mid
        \text{if $q$ is injective, then for every $r$ such that
            $r$ is a finite sequence in $A$ and
            $r$ is injective and $\ran{q}=\ran{r}$ we have
            $\sequenceFold{\mul[M]}{\neutral[M]}{f\circ q}
                =\sequenceFold{\mul[M]}{\neutral[M]}{f\circ r}$}\}$.
    Show that $\emptyset\in C$.
    \begin{subproof}
        $\emptyset$ is a finite sequence in $A$ by
            \auto{emptyset_finite_sequence}.
        Hence $\emptyset\in\finiteSequences{A}$ by
            \auto{finite_sequences_iff}.
        $\emptyset$ is injective by \auto{emptyset,injective}.
        Show that for every $r$ such that
            $r$ is a finite sequence in $A$ and
            $r$ is injective and $\ran{\emptyset}=\ran{r}$ we have
            $\sequenceFold{\mul[M]}{\neutral[M]}{f\circ\emptyset}
                =\sequenceFold{\mul[M]}{\neutral[M]}{f\circ r}$.
        \begin{subproof}
            Fix $r$.
            Assume $r$ is a finite sequence in $A$ and
                $r$ is injective and $\ran{\emptyset}=\ran{r}$.
            Hence $\ran{r}=\emptyset$ by \auto{ran_emptyset}.
            $r$ is a relation by
                \auto{finite_sequence_domain,funs_is_relation}.
            Hence $r=\emptyset$ by
                \auto{setext,emptyset,relation,ran_intro}.
            Follows by \auto.
        \end{subproof}
        Follows by \auto.
    \end{subproof}
    Show that for every $q,y$ such that
        $q\in C$ and $y\in A$ and $q$ is a finite sequence in $A$
        we have $\sequenceAppend{q}{y}\in C$.
    \begin{subproof}
        Fix $q,y$.
        Assume $q\in C$ and $y\in A$ and
            $q$ is a finite sequence in $A$.
        Let $w=\sequenceAppend{q}{y}$.
        $w$ is a finite sequence in $A$ by
            \auto{sequence_append_finite_sequence}.
        Hence $w\in\finiteSequences{A}$ by
            \auto{finite_sequences_iff}.
        Show that if $w$ is injective, then for every $r$ such that
            $r$ is a finite sequence in $A$ and
            $r$ is injective and $\ran{w}=\ran{r}$ we have
            $\sequenceFold{\mul[M]}{\neutral[M]}{f\circ w}
                =\sequenceFold{\mul[M]}{\neutral[M]}{f\circ r}$.
        \begin{subproof}
            Assume $w$ is injective.
            Hence $q$ is injective and $y\notin\ran{q}$ by
                \auto{spike_sequence_append_injective_elim}.
            Fix $r$.
            Assume $r$ is a finite sequence in $A$ and
                $r$ is injective and $\ran{w}=\ran{r}$.
            $\ran{w}=\cons{y}{\ran{q}}$ by
                \auto{spike_sequence_append_range}.
            Hence $y\in\ran{r}$ by \auto{cons_iff}.
            Take $v$ such that
                $v$ is a finite sequence in $A$ and
                $v$ is injective and
                $\ran{v}=\ran{r}\setminus\{y\}$ and
                $\sequenceFold{\mul[M]}{\neutral[M]}{f\circ r}
                    =\mul[M]
                        (\sequenceFold{\mul[M]}{\neutral[M]}{f\circ v},
                         f(y))$ by
                \auto{spike_fold_extract_value}.
            $\ran{q}=\ran{r}\setminus\{y\}$ by
                \auto{setext,setminus,cons_iff,singleton_iff}.
            Hence $\ran{q}=\ran{v}$ by \auto.
            $\sequenceFold{\mul[M]}{\neutral[M]}{f\circ q}
                =\sequenceFold{\mul[M]}{\neutral[M]}{f\circ v}$ by
                \auto.
            $\sequenceFold{\mul[M]}{\neutral[M]}{f\circ w}
                =\mul[M]
                    (\sequenceFold{\mul[M]}{\neutral[M]}{f\circ q},
                     f(y))$ by
                \auto{spike_mapped_append_fold}.
            Follows by \auto.
        \end{subproof}
        Follows by \auto.
    \end{subproof}
    $s\in C$ by \auto{finite_sequence_induction}.
    Follows by \auto.
\end{proof}

\begin{corollary}\label{spike_enumeration_fold_invariant}
    Let $M$ be a monoid.
    Suppose for all $u,v\in\carrier[M]$ we have
        $\mul[M](u,v)=\mul[M](v,u)$.
    Suppose $e$ is an enumeration of $A$ and
        $d$ is an enumeration of $A$.
    Suppose $f\in\funs{A}{\carrier[M]}$.
    Then
        $\sequenceFold{\mul[M]}{\neutral[M]}{f\circ e}
            =\sequenceFold{\mul[M]}{\neutral[M]}{f\circ d}$.
\end{corollary}
\begin{proof}
    $e$ is a finite sequence in $A$ and
        $d$ is a finite sequence in $A$ by
        \auto{enumeration_is_finite_sequence}.
    $e$ is injective and $d$ is injective and
        $\ran{e}=A=\ran{d}$ by
        \auto{enumeration,bijections,ran_of_surj}.
    Follows by \auto{spike_injective_sequence_fold_invariant}.
\end{proof}

\begin{definition}\label{spike_finite_fold_values}
    $\spikeFiniteFoldValues{M}{A}{f}=
        \{y\in\carrier[M]\mid
            \text{there exists $e$ such that
                $e$ is an enumeration of $A$ and
                $y=\sequenceFold{\mul[M]}{\neutral[M]}{f\circ e}$}\}$.
\end{definition}

\begin{proposition}\label{spike_finite_fold_values_iff}
    $y\in\spikeFiniteFoldValues{M}{A}{f}$ iff
        $y\in\carrier[M]$ and there exists $e$ such that
            $e$ is an enumeration of $A$ and
            $y=\sequenceFold{\mul[M]}{\neutral[M]}{f\circ e}$.
\end{proposition}
\begin{proof}
    Follows by \auto{spike_finite_fold_values}.
\end{proof}

\begin{theorem}\label{spike_finite_fold_values_inhabited_subsingleton}
    Let $M$ be a monoid.
    Suppose for all $u,v\in\carrier[M]$ we have
        $\mul[M](u,v)=\mul[M](v,u)$.
    Suppose $A$ is finite and $f\in\funs{A}{\carrier[M]}$.
    Then $\spikeFiniteFoldValues{M}{A}{f}$ is inhabited and
        $\spikeFiniteFoldValues{M}{A}{f}$ is a subsingleton.
\end{theorem}
\begin{proof}
    Take $e$ such that $e$ is an enumeration of $A$ by
        \auto{finite_iff_has_enumeration}.
    $e$ is a finite sequence in $A$ by
        \auto{enumeration_is_finite_sequence}.
    Let $p=\sequenceFold{\mul[M]}{\neutral[M]}{f\circ e}$.
    $p\in\carrier[M]$ by \auto{spike_mapped_fold_type}.
    Hence $p\in\spikeFiniteFoldValues{M}{A}{f}$ by
        \auto{spike_finite_fold_values_iff}.
    Hence $\spikeFiniteFoldValues{M}{A}{f}$ is inhabited by \auto.
    Show that $\spikeFiniteFoldValues{M}{A}{f}$ is a subsingleton.
    \begin{subproof}
        Show that for all $u,v\in\spikeFiniteFoldValues{M}{A}{f}$
            we have $u=v$.
        \begin{subproof}
            Fix $u,v\in\spikeFiniteFoldValues{M}{A}{f}$.
            Take $d$ such that $d$ is an enumeration of $A$ and
                $u=\sequenceFold{\mul[M]}{\neutral[M]}{f\circ d}$ by
                \auto{spike_finite_fold_values_iff}.
            Take $q$ such that $q$ is an enumeration of $A$ and
                $v=\sequenceFold{\mul[M]}{\neutral[M]}{f\circ q}$ by
                \auto{spike_finite_fold_values_iff}.
            Follows by \auto{spike_enumeration_fold_invariant}.
        \end{subproof}
        Follows.
    \end{subproof}
    Follows.
\end{proof}

\begin{theorem}\label{spike_finite_fold_values_singleton}
    Let $M$ be a monoid.
    Suppose for all $u,v\in\carrier[M]$ we have
        $\mul[M](u,v)=\mul[M](v,u)$.
    Suppose $A$ is finite and $f\in\funs{A}{\carrier[M]}$.
    Then there exists $p\in\carrier[M]$ such that
        $\spikeFiniteFoldValues{M}{A}{f}=\{p\}$.
\end{theorem}
\begin{proof}
    $\spikeFiniteFoldValues{M}{A}{f}$ is inhabited and
        $\spikeFiniteFoldValues{M}{A}{f}$ is a subsingleton by
        \auto{spike_finite_fold_values_inhabited_subsingleton}.
    Take $p$ such that
        $p\in\spikeFiniteFoldValues{M}{A}{f}$ by \auto.
    Hence $p\in\carrier[M]$ by
        \auto{spike_finite_fold_values_iff}.
    $\spikeFiniteFoldValues{M}{A}{f}=\{p\}$ by
        \auto{singleton_iff_inhabited_subsingleton}.
    Follows by \auto.
\end{proof}

\begin{proposition}\label{spike_unions_singleton}
    $\unions{\{x\}}=x$.
\end{proposition}
\begin{proof}
    Follows by \auto{setext,unions_iff,singleton_iff}.
\end{proof}

\begin{definition}\label{spike_finite_fold}
    $\finiteFold{M}{A}{f}=\unions{\spikeFiniteFoldValues{M}{A}{f}}$.
\end{definition}

\begin{theorem}\label{spike_finite_fold_canonical}
    Let $M$ be a monoid.
    Suppose for all $u,v\in\carrier[M]$ we have
        $\mul[M](u,v)=\mul[M](v,u)$.
    Suppose $A$ is finite and $f\in\funs{A}{\carrier[M]}$.
    Then
        $\finiteFold{M}{A}{f}\in\carrier[M]$ and
        $\spikeFiniteFoldValues{M}{A}{f}
            =\{\finiteFold{M}{A}{f}\}$ and
        for every $e$ such that $e$ is an enumeration of $A$ we have
            $\finiteFold{M}{A}{f}
                =\sequenceFold{\mul[M]}{\neutral[M]}{f\circ e}$.
\end{theorem}
\begin{proof}
    Take $p\in\carrier[M]$ such that
        $\spikeFiniteFoldValues{M}{A}{f}=\{p\}$ by
        \auto{spike_finite_fold_values_singleton}.
    $\finiteFold{M}{A}{f}=\unions{\{p\}}$ by
        \auto{spike_finite_fold}.
    Hence $\finiteFold{M}{A}{f}=p$ by
        \auto{spike_unions_singleton}.
    Hence $\finiteFold{M}{A}{f}\in\carrier[M]$ and
        $\spikeFiniteFoldValues{M}{A}{f}
            =\{\finiteFold{M}{A}{f}\}$ by \auto.
    Show that for every $e$ such that $e$ is an enumeration of $A$
        we have
            $\finiteFold{M}{A}{f}
                =\sequenceFold{\mul[M]}{\neutral[M]}{f\circ e}$.
    \begin{subproof}
        Fix $e$.
        Assume $e$ is an enumeration of $A$.
        $e$ is a finite sequence in $A$ by
            \auto{enumeration_is_finite_sequence}.
        Hence
            $\sequenceFold{\mul[M]}{\neutral[M]}{f\circ e}
                \in\carrier[M]$ by \auto{spike_mapped_fold_type}.
        Hence
            $\sequenceFold{\mul[M]}{\neutral[M]}{f\circ e}
                \in\spikeFiniteFoldValues{M}{A}{f}$ by
            \auto{spike_finite_fold_values_iff}.
        Follows by \auto{singleton_iff}.
    \end{subproof}
    Follows.
\end{proof}
